\documentclass{beamer}

\newenvironment{tightcenter}{%
  \setlength\topsep{0pt}
  \setlength\parskip{0pt}
  \begin{center}
}{%
  \end{center}
}

\newenvironment{Snippet}{\Verbatim[samepage=true,fontsize=\tiny]}{\endVerbatim}

\mode<presentation>
{
  \usetheme{Copenhagen}
  %%\usecolortheme[RGB={173,222,25}]{structure}
  % \usecolortheme[RGB={255,0,0}]{structure}
  \setbeamertemplate{items}[circle]
  \setbeamercovered{transparent}
}

\usepackage[polish]{babel}
\usepackage{chessfss}
\usepackage{hyperref}
\usepackage{qtree}
\usepackage{mathtools}
\usepackage{dirtytalk}
\usepackage{epigraph}
\usepackage{textgreek}
\usepackage[utf8]{inputenc}
\usepackage{times}
\usepackage[T1]{fontenc}
\usepackage{tikz}
\usepackage{csquotes}
\usepackage{amsmath}
\usepackage{fancyvrb}
\usepackage{ulem}
\usepackage{adjustbox}


\title{\textbf{Uogólnione Programowanie Wizualne}}

\author{mgr Maciej Godek}

\institute{
  \tiny{\href{mailto:m.godek@wi.umg.edu.pl}{\textbf{m.godek@wi.umg.edu.pl}}}
}

\date{\textbf{Uniwersytet Morski}, 06.03.2026}


\begin{document}

\begin{frame}
  \titlepage
\end{frame}

\begin{frame}
  Czym jest \textit{Uogólnione Programowanie Wizualne}?
  \begin{itemize}
    \pause \item krótka odpowiedź: nie ma czegoś takiego
    \pause \item nieco dłuższa odpowiedź: coś co sam sobie wymyśliłem
  \end{itemize}
\end{frame}


\begin{frame}[fragile]
  Motywacja: przykłady wydają się prostsze do zrozumienia niż definicje.

  \begin{columns}[T]
    \begin{column}{.58\textwidth}
      \only<1->{
      \begin{Snippet}
(define (multicombinations A n)
  (cond
    ((= n 0)
     '())
    ((= n 1)
     (map list A))
    (else
     (append-map
       (lambda (combination)
         (map (lambda (a)
                `(,a . ,combination))
              A))
       (multicombinations A (- n 1))))))
      \end{Snippet}
      }
    \end{column}

    \begin{column}{.38\textwidth}
      \only<2->{
      \begin{Snippet}
(e.g.
 (multicombinations '(a b) 3)
 ===>
 ((a a a)
  (b a a)
  (a b a)
  (b b a)
  (a a b)
  (b a b)
  (a b b)
  (b b b)))
      \end{Snippet}
      }
    \end{column}
  \end{columns}
\end{frame}



\end{document}
